\documentclass[11pt,a4paper]{article}

% ===================== Codificación y fuentes =====================
\usepackage[utf8]{inputenc}
\usepackage[T1]{fontenc}
\usepackage{lmodern}

% ===================== Página y utilidades ========================
\usepackage[left=2.5cm,right=2.5cm,top=2.5cm,bottom=2.5cm]{geometry}
\usepackage{array}
\usepackage{enumitem}
\usepackage{graphicx}
\usepackage{xcolor}
\usepackage{titlesec}
\usepackage{fancyhdr}
\usepackage[most]{tcolorbox}
\usepackage{booktabs}
\usepackage{longtable}
\usepackage{float}
\usepackage{amsmath}
\usepackage{caption}
\usepackage{microtype}
\usepackage{hyperref}

% ===================== Etiquetas en español (sin babel) ===========
\renewcommand{\contentsname}{Índice}
\renewcommand{\figurename}{Figura}
\renewcommand{\tablename}{Tabla}
\renewcommand{\refname}{Bibliografía}

% ===================== Escala de grises (sin color) ================
\definecolor{negro}{RGB}{0,0,0}
\definecolor{gris}{RGB}{90,90,90}
\definecolor{grisclaro}{RGB}{150,150,150}
\definecolor{fillgray}{RGB}{175,175,175}

% ===================== Hipervínculos (negro) =======================
\hypersetup{colorlinks=true, linkcolor=negro, urlcolor=negro, citecolor=negro,
  pdftitle={Evaluación Sumativa 3: Cierre y comunicación del proyecto},
  pdfauthor={José Ignacio Rodríguez, Yerko Gallardo Miranda, Sebastián Rojas Araya}}

% ===================== Estilo de títulos ===========================
\titleformat{\section}{\large\bfseries\color{negro}}{\thesection.}{0.5em}{}
\titleformat{\subsection}{\normalsize\bfseries\color{negro}}{\thesubsection}{0.5em}{}
\titlespacing*{\section}{0pt}{1.0em}{0.3em}
\titlespacing*{\subsection}{0pt}{0.7em}{0.2em}

% ===================== Pie de página ===============================
\pagestyle{fancy}
\fancyhf{}
\fancyfoot[L]{\footnotesize\color{gris}Magíster en Ciencia de Datos e Inteligencia Artificial}
\fancyfoot[R]{\footnotesize\color{gris}\thepage}
\renewcommand{\headrulewidth}{0pt}
\renewcommand{\footrulewidth}{0.4pt}

% ===================== Ajustes del informe =========================
\captionsetup{font=small,labelfont=bf,skip=3pt}
\setlength{\floatsep}{6pt}\setlength{\textfloatsep}{8pt}\setlength{\intextsep}{6pt}
\setlist{nosep,leftmargin=1.2em}
\newcommand{\fig}[4]{\begin{figure}[!htb]\centering\includegraphics[width=#1\linewidth]{../../outputs/figuras/#2}\caption{#3}\label{#4}\end{figure}}
\newcommand{\codigo}[1]{{\small\texttt{#1}}}

% ===================== Datos del grupo ==============================
\newcommand{\proyTitulo}{Factores asociados al atraso en envíos de comercio electrónico: del análisis exploratorio a la clasificación}
\newcommand{\proyIntegrantes}{José Ignacio Rodríguez, Yerko Gallardo Miranda y Sebastián Rojas Araya}
\newcommand{\proyDataset}{\emph{E-Commerce Shipping Data} (Gopalani, 2021), 10.999 envíos}
\newcommand{\proyDocente}{J. P. Maidana}
\newcommand{\proyFecha}{28/10/2026}
\newcommand{\proyRepo}{\url{https://github.com/joserodriguezc/grupo_1_estadistica_computacional}}

\setlength{\parindent}{0pt}
\setlength{\parskip}{0.3em}

\begin{document}
\pagenumbering{gobble}

\thispagestyle{empty}
\begin{center}
\vspace*{1.0cm}
\includegraphics[height=4.4cm]{logo_unab.png}\\[1.3cm]

{\footnotesize\color{gris}\MakeUppercase{Magíster en Ciencia de Datos e Inteligencia Artificial}}\\[0.15cm]
{\footnotesize\color{grisclaro}MCDI501: Estadística Computacional para la Toma de Decisiones}\\[2.0cm]

{\Large\bfseries\color{negro} Evaluación Sumativa 3}\\[0.45cm]
{\large\color{negro} Cierre y Comunicación del Proyecto}\\[1.0cm]

{\small\itshape\color{gris} Informe final integrado \textperiodcentered\ 60\,\%}
\end{center}

\vspace{1.6cm}

\begin{flushleft}
\setlength{\extrarowheight}{5pt}
\begin{tabular}{@{}>{\bfseries\color{negro}}l >{\raggedright\arraybackslash}p{10.5cm}@{}}
Título del proyecto: & \proyTitulo \\
Integrantes:         & \proyIntegrantes \\
Dataset:             & \proyDataset \\
Docente:             & \proyDocente \\
Fecha:               & \proyFecha \\
Repositorio GitHub:  & \proyRepo \\
\end{tabular}
\end{flushleft}

\clearpage
\pagenumbering{arabic}

% =====================================================================
% Máximo 10 páginas (portada aparte). Presupuesto y responsables según la planificación.
% Cada sección declara qué fase (S1, S2, S3) informa a la siguiente.
% Todas las cifras deben provenir de sumativa_3/outputs/tablas/.
% =====================================================================

\section*{Resumen ejecutivo}
% 0,5 páginas - José. Problema, enfoque S1 -> S2 -> S3, resultado principal y recomendación accionable.
\emph{(Pendiente.)}

\section{Introducción}
% 0,75 páginas - José. Contexto, pregunta de S1, qué aportó cada fase y objetivo de S3.
\emph{(Pendiente.)}

\section{Metodología integrada}
% 2,5 páginas en total.
\subsection{Datos, trazabilidad y punto de partida}
% José. Datos de S1 con huella, semilla del beacon, partición común; hallazgos de S1 y S2 que guían S3 (figura 0).
\emph{(Pendiente.)}
\subsection{Datos faltantes e imputación por regresión}
% Yerko. Faltantes en S1, simulación MAR, predictores (S1 + estabilidad S2), supuestos y estrategias.
\emph{(Pendiente.)}
\subsection{Regresión logística: preparación, selección y evaluación}
% Sebastián. Outliers, codificación, M1 (S1 + S2), M2 (stepwise), M3 (AIC/BIC), métricas y criterio de selección.
\emph{(Pendiente.)}
\subsection{Estabilidad y diagnóstico}
% Yerko (bootstrap) y Sebastián (VIF, linealidad en el logit, influyentes y residuos).
\emph{(Pendiente.)}

\section{Resultados integrados}
% 4 páginas en total.
\subsection{Datos faltantes e imputación}
% Yerko.
\emph{(Pendiente.)}
\subsection{Modelos logísticos y desempeño}
% Sebastián.
\emph{(Pendiente.)}
\subsection{Estabilidad del modelo final}
% Yerko.
\emph{(Pendiente.)}
\subsection{Diagnóstico del modelo final}
% Sebastián.
\emph{(Pendiente.)}
\subsection{Impacto de la estrategia de imputación}
% José.
\emph{(Pendiente.)}

\section{Discusión}
% 1,25 páginas - los tres; José integra. Qué se confirma de S1 y S2; límites (descuento, bloque 1-3.135,
% MAR simulado, causalidad) y validez externa.
\emph{(Pendiente.)}

\section{Conclusiones y recomendaciones}
% 0,75 páginas - José. Respuestas accionables para la empresa y para un análisis futuro.
\emph{(Pendiente.)}

\section*{Referencias}
% 0,25 páginas - José.
\begin{small}
\begin{itemize}
\item Efron, B. y Tibshirani, R. J. (1993). \emph{An Introduction to the Bootstrap}. Chapman \& Hall.
\item Gopalani, P. (2021). \emph{E-Commerce Shipping Data}, versión 1. Kaggle.
\item Hosmer, D. W., Lemeshow, S. y Sturdivant, R. X. (2013). \emph{Applied Logistic Regression} (3.\textsuperscript{a} ed.). Wiley.
\item Little, R. J. A. y Rubin, D. B. (2019). \emph{Statistical Analysis with Missing Data} (3.\textsuperscript{a} ed.). Wiley.
\end{itemize}
\end{small}

\end{document}
